\section{Planteamiento del problema}
\label{sec:planteamiento}

Hacen falta tres cortes. Si los tres cuentan el TSB, sobran dos.

\subsection{Situación real (problemática)}

Los casos y las cifras quedan en §\ref{sec:problematica}: el TSB, la orden de la CMA, el 40\,\%
de Casas et al.\ (2021) y la revisión de Golmohammadi et al.\ (2024). Ahí hay clientes fuera de
servicio y APIs exigidas por regulación. No hay un contrato elegido ni una fila de la matriz.

\subsection{Situación subyacente (problema abstracto)}

Sin el nombre del banco quedan dos estructuras que no coinciden de antemano. El contrato es un
árbol de \textit{paths}, operaciones y \textit{schemas}. El \ServiceDomain{} es otro:
\Behavior{}s, \BusinessObject{}s y relaciones de control. El problema abstracto es esa
correspondencia. Los tipos ---nombre, verbo HTTP, objeto no cubierto, versión del YAML que no
revisa el modelo--- están en la \Cref{tab:tipologias}. El par \texttt{PaymentOrder} /
\texttt{paymentRequest} ya se usó en §\ref{sec:problematica}; aquí no se repite.

De ese corte salen el objeto y la unidad de §\ref{sec:objeto-estudio}: el contrato completo, y
cada par endpoint/schema $\leftrightarrow$ \Behavior{}/\BusinessObject{}. BIAN entra como
documento de referencia, no como el \textit{core} del banco.

\subsection{Problema tecnológico (situación a modificar)}

Lo que un equipo tiene hoy es el Swagger leído a ojo, o un \textit{schema matcher} que no distingue
\ServiceDomain{}s. Para el par contrato $\times$ SD no hay un \StructuralScore{}, un
\SemanticScore{}, un \CoverageScore{} ni un nivel (Alta, Media, Baja, Nula) con los umbrales
escritos. El objetivo general del Capítulo~\ref{cap:objetivos} y la matriz del Anexo~C describen
ese procedimiento. Correrlo sobre el \textit{core} de un banco queda fuera de esta Guía
(\Cref{sec:alcances}).

\subsection{Problema de investigación}

El \textbf{problema de investigación} es una pregunta, no un objetivo:

\begin{quote}
  \textbf{¿En qué medida, y con qué tipos de desajuste, un contrato OpenAPI bancario se aparta
  del \ServiceDomain{} BIAN emparejado, y cómo se convierte eso en un score?}
\end{quote}

El §\ref{sec:brecha} dice qué parte de esa pregunta dejan abierta los 14 papers.

Las \textbf{preguntas específicas} bajan el mismo hueco a cuatro cortes, no a cuatro objetivos:

\begin{table}[H]
  \centering
  \caption{Preguntas específicas de investigación}
  \label{tab:preguntas}
  \small
  \begin{tabular}{@{}p{0.06\textwidth}p{0.88\textwidth}@{}}
    \toprule
    \textbf{ID} & \textbf{Pregunta} \\
    \midrule
    \textbf{P1} &
    ¿Qué correspondencias estructurales pueden establecerse entre elementos de un contrato OpenAPI
    (\textit{paths}, \textit{operations}, \textit{schemas}) y los artefactos de un \ServiceDomain{}
    BIAN? \\
    \textbf{P2} &
    ¿Qué grado de similitud semántica existe entre los términos y conceptos de negocio en BIAN y los
    nombres de recursos, propiedades y descripciones en OpenAPI? \\
    \textbf{P3} &
    ¿Qué tipologías de desalineación predominan al comparar ambos artefactos en un mismo dominio de
    servicio? \\
    \textbf{P4} &
    ¿Es posible integrar métricas estructurales y semánticas en un indicador de coherencia global
    por cada instancia evaluada (par contrato OpenAPI bancario --- \ServiceDomain{} BIAN)? \\
    \bottomrule
  \end{tabular}
\end{table}

La \Cref{tab:tipologias} resume la tipología de desalineaciones identificadas:

\begin{table}[H]
  \centering
  \caption{Tipología de desalineaciones OpenAPI--BIAN}
  \label{tab:tipologias}
  \begin{tabular}{@{}p{0.18\textwidth}p{0.38\textwidth}p{0.32\textwidth}@{}}
    \toprule
    \textbf{Tipo} & \textbf{Descripción} & \textbf{Ejemplo ilustrativo} \\
    \midrule
    \textbf{Estructural} &
    Divergencia entre la organización de endpoints/esquemas y la estructura del SD BIAN &
    Operación \texttt{POST /transfer} sin correspondencia con el \Behavior{} \texttt{Initiate} definido en BIAN \\
    \textbf{Semántica} &
    Equivalencia estructural parcial pero significado de negocio distinto &
    Propiedad \texttt{amount} sin vínculo explícito con \texttt{PaymentAmount} del modelo BIAN \\
    \textbf{De modelado} &
    Representación incorrecta o incompleta de entidades de negocio &
    Esquema que omite atributos obligatorios del \BusinessObject{} BIAN \\
    \textbf{REST / contrato} &
    Inconsistencias en el estilo REST que dificultan la interpretación alineada al dominio &
    Uso de verbos en URLs o recursos que no mapean a \textit{Service Operations} BIAN \\
    \bottomrule
  \end{tabular}
\end{table}


P1--P4 no son los objetivos. Normalizar, emparejar y puntuar están en OE1--OE4
(Capítulo~\ref{cap:objetivos}) y en las cinco fases del Anexo~C.
